\subsection{APPELL POLYNOMIALS ATTACHED TO A COMPOSITION OPERATOR}

Given a composition operator \(U = \sum_{k=0}^{\infty} \alpha_k \mathrm{D}^k \neq 0\) let \(p\) be the least of the integers \(k\) such that \(\alpha_k \neq 0\) ; we shall say that \(p\) is the order of the operator \(U\) .

PROPOSITION 2. Every composition operator of order 0 is invertible in the algebra \(\Gamma\) of composition operators on K{[}X{]}.

Indeed, a formal series \(\sum_{k=0}^{\infty}\alpha_{k}S^{k}\) such that \(\alpha_{0}\neq0\) is invertible in the ring K{[}{[}S{]}{]} (Alg., IV. 41); the proposition thus follows from th. 1 of VI, p.~271.

PROPOSITION 3. Let U be a composition operator of order p; then \(U(f) = 0\) for every polynomial f of degree \textless{} p; for every polynomial \(f \neq 0\) of degree \(n \geqslant p\) , \(U(f)\) is a polynomial \(\neq 0\) of degree n - p.

This is an immediate consequence of formula (2) of VI, p.~271 and of the definition of the order of \(U\) .

It is clear that every operator \(U\) of order \(p\) can be written in a unique way as \(U = \mathrm{D}^p V = V\mathrm{D}^p\) , where \(V\) is an operator of order 0 (and so invertible).

DEFINITION 2. Let \(U = \mathrm{D}^p V\) be a composition operator of order \(p\) on \(\mathrm{K}[\mathrm{X}]\) . The polynomial \(u_n(\mathrm{X}) = V^{-1}(\mathrm{X}^n)\) is called the Appell polynomial of index \(n\) attached to the operator \(U\) .

\[
\text {   If   } V ^ {- 1} = \sum_ {k = 0} ^ {\infty} \frac {1}{k !} \beta_ {k} D ^ {k} (\text {   with   } \beta_ {0} \neq 0) \text {   one   thus   has   }
\]

\[
u _ {n} (\mathrm{X}) = \sum_ {k = 0} ^ {n} {\binom {n} {k}} \beta_ {k}   \mathrm{X} ^ {n - k}.\tag{6}
\]

One verifies that \(u_{n}\) is a polynomial of degree \(n\) (prop. 3); further

\[
u _ {n} (0) = \beta_ {n}.
\]

PROPOSITION 4. The Appell polynomials attached to \(U\) satisfy the relations

\[
\frac {d u _ {n}}{d \mathrm{X}} = n u _ {n - 1}\tag{7}
\]

\[
u _ {n} (\mathrm{X} + \mathrm{Y}) = \sum_ {k = 0} ^ {n} {\binom {n} {k}} u _ {n - k} (\mathrm{X})   \mathrm{Y} ^ {k}\tag{8}
\]

\[
U \left(u _ {n} (\mathrm{X})\right) = \frac {n !}{(n - p) !} \mathrm{X} ^ {n - p}.\tag{9}
\]

These formulae are in fact respectively equivalent to the following relations (bearing def. 2 in mind):

\[
\mathbf {D} V ^ {- 1} = V ^ {- 1} \mathbf {D}\tag{10}
\]

\[
\left(\exp \left(\mathrm{YD} _ {\mathrm{X}}\right)\right) V _ {\mathrm{X}} ^ {- 1} = V _ {\mathrm{X}} ^ {- 1} \exp \left(\mathrm{YD} _ {\mathrm{X}}\right)\tag{11}
\]

\[
U V ^ {- 1} = \mathbf {D} ^ {p}.\tag{12}
\]

PROPOSITION 5. For every polynomial \(f \in \mathrm{K}[\mathrm{X}]\) and every composition operator \(U\) of order \(p\) one has

\[
f ^ {(p)} (\mathrm{X} + \mathrm{Y}) = \sum_ {k = 0} ^ {\infty} \frac {1}{k !} U (f ^ {(k)} (\mathrm{X})) u _ {k} (\mathrm{Y})\tag{13}
\]

(generalized Taylor expansion).

Indeed, if one puts \(U = \mathrm{D}^p V = V\mathrm{D}^p\) one has (VI, p.~270, formula (1))

\[
V _ {\mathrm{X}} ^ {- 1} (f (\mathrm{X} + \mathrm{Y})) = V _ {\mathrm{Y}} ^ {- 1} (f (\mathrm{X} + \mathrm{Y})) = \sum_ {k = 0} ^ {\infty} \frac {1}{k !} f ^ {(k)} (\mathrm{X}) u _ {k} (\mathrm{Y})\tag{14}
\]

by the Taylor formula and by def. 2 of VI, p.~273; it suffices to apply the operator \(U_{\mathrm{X}}\) to the first and last terms of formula (14) to obtain (13).

